\chapter{ЗАГРУЗКА И ОТРИСОВКА МОДЕЛЕЙ РОБОТОВ В ФОРМАТЕ URDF}
\label{cha:urdf}

\section{Постановка задачи}

URDF (Unified Robot Description Format)~\cite{urdf} --- основанный на XML формат
описания кинематической структуры робота: набор твёрдых звеньев (\texttt{link}),
соединённых суставами (\texttt{joint}). Каждое звено содержит визуальную
(\texttt{visual}), столкновительную (\texttt{collision}) и инерционную
(\texttt{inertial}) геометрию; сустав задаёт тип сочленения (вращательное,
призматическое, фиксированное и др.), оси, пределы и взаимное расположение
родительского и дочернего звеньев. Задача состоит в загрузке URDF-модели,
построении дерева звеньев и его визуализации в трёхмерной сцене с возможностью
интерактивного управления.

Работа выполнена в репозитории \texttt{modeuler}, модуль \texttt{scene}.

\section{Алгоритм загрузки}

Разбор URDF выполняется в два прохода. На первом проходе извлекаются все звенья и
суставы с их атрибутами и геометрией. На втором --- по родительско-дочерним
связям суставов строится дерево звеньев, вычисляются относительные
преобразования и определяется корневое звено (звено, не являющееся дочерним ни
для одного сустава). Преобразования заданы в виде однородных матриц $4\times4$;
положение звена в мировой системе координат получается перемножением
преобразований вдоль пути от корня.

\begin{lstlisting}[language=C++,caption={Разбор URDF и построение дерева звеньев (схематично)},label=lst:urdf]
UrdfModel load_urdf(const std::filesystem::path& path) {
    auto xml = parse_xml(path);                    // pugixml
    UrdfModel model;
    for (auto& link : xml.select("robot/link"))
        model.links.emplace(link.attr("name"), parse_link(link));
    for (auto& joint : xml.select("robot/joint")) {
        Joint j = parse_joint(joint);              // type, axis, origin, limits
        model.links[j.child].parent_joint = j.name;
        model.joints.emplace(j.name, std::move(j));
    }
    model.root = find_root_link(model);            // звено без родительского сустава
    compute_world_transforms(model, model.root);   // обход дерева в глубину
    return model;
}
\end{lstlisting}

Каждому визуальному элементу звена сопоставляется примитив (бокс, цилиндр, сфера)
или загружаемая полигональная сетка (mesh). Загрузка сеток и материалов
выполняется средствами библиотеки vsgXchange, что обеспечивает поддержку
распространённых форматов геометрии. Описание звеньев и суставов хранится в
структурах \texttt{Link} и \texttt{Joint} (заголовки \texttt{models/link.hpp},
\texttt{models/joint.hpp}).

\section{Отрисовка}

Визуализация построена на VulkanSceneGraph (VSG)~\cite{vsg} --- современной
библиотеке графа сцены поверх графического API Vulkan~\cite{vulkan}. Для каждого звена создаётся узел графа
сцены с матрицей преобразования; визуальная геометрия звена присоединяется как
дочерний узел. Это позволяет при изменении обобщённых координат сустава
обновлять лишь соответствующую матрицу, не пересобирая граф сцены.

Подсистема визуализации включает (модуль \texttt{scene}):
\begin{itemize}
  \item управление камерой и навигацией (\texttt{utils/camera.hpp},
        \texttt{controllers/arcball\_controller}) --- орбитальное вращение,
        панорамирование, масштабирование;
  \item выбор объектов сцены (\texttt{controllers/picking\_controller}) ---
        определение звена под курсором методом трассировки луча;
  \item интерактивное вращение суставов
        (\texttt{controllers/\allowbreak joint\_rotation\_controller}) --- изменение
        обобщённой координаты вращательного сустава перетаскиванием;
  \item вспомогательные элементы (\texttt{utils/arrow\_gizmo}) --- визуализация
        осей и манипуляторов.
\end{itemize}

\screenshot{0.85}{Загруженная URDF-модель робота в трёхмерной сцене}{Снимок
экрана: окно просмотра \texttt{modeuler} с загруженной моделью манипулятора в
формате URDF; видны звенья, оси суставов и элементы управления камерой.}{urdf-view}

\section{Результаты}

Реализованы загрузка URDF-моделей, построение дерева звеньев с вычислением
мировых преобразований и интерактивная визуализация в реальном времени.
Поддержаны основные типы суставов и базовые визуальные примитивы, а также
загрузка полигональных сеток. Корректность построения дерева и преобразований
проверена на тестовых моделях с известной кинематикой; результат соответствует
ожидаемому расположению звеньев.
